\chapter{Consumer Groups, Rebalancing, and Commit Fencing}
A consumer group is a set of consumers cooperating to read one topic: a member is a participant identifier, an assignment is the current list of partitions owned by each member, and one partition is assigned to at most one member in the group at a time. The coordinator stores group state and answers join/assignment requests; membership changes make it recompute the whole assignment, a change called a rebalance. A generation is the group version incremented on each topology change, not a clock or consumer offset. Later, a GroupToken carries group, member, and generation together so the server can reject stale members. Heartbeats indicate that a member is still active; expired members are removed. The course models one topic, immediate rebalancing, and a trusted in-process coordinator, not the full Kafka group protocol or highly available consensus.

\begin{figure}[ht]
\centering
\begin{tikzpicture}[font=\small,node distance=10mm]
\node[draw,rounded corners,minimum width=19mm] (a1) {A joins};
\node[draw,rounded corners,minimum width=26mm,right=of a1] (g1) {generation 1};
\node[draw,rounded corners,minimum width=19mm,right=of g1] (b) {B joins};
\node[draw,rounded corners,minimum width=26mm,right=of b] (g2) {generation 2};
\node[draw,rounded corners,minimum width=23mm,below=of g2] (fence) {A gen 1 fenced};
\draw[-{Stealth},thick] (a1)--(g1); \draw[-{Stealth},thick] (g1)--(b);
\draw[-{Stealth},thick] (b)--(g2); \draw[-{Stealth},thick] (g2)--(fence);
\end{tikzpicture}
\caption{A membership topology change increments the generation; an old token cannot overwrite the current commit.}
\end{figure}

\lessonsection{17}{Partition Assignment}
\goals{Implement deterministic and balanced Round Robin assignment with exactly one owner per partition.}{Understand deduplication, stable sorting, immutable results, and quotient/remainder distribution.}
\background{Partitions can be processed in parallel by consumers, but two members processing one partition at the same time can conflict over progress. Round Robin cycles through sorted members to assign sorted partitions; its result is repeatable and partition counts differ by at most one. This assignor is a pure calculation and does not determine whether a member has actually joined a group.}
\subsubsection*{Read first}
Review \stepref{3}{TopicPartition and partition ordering}. The assignor is called by the \stepref{18}{membership coordinator}; it allocates only the partition/member lists it receives and does not query a catalog or modify a log.
\providededit{The \method{TopicPartition} value type and group result type are provided; the \method{RoundRobinAssignor} constructor is stateless and does not provide the assignment algorithm.}{\pathcode{exercises/src/main/java/io/simplekafka/group/RoundRobinAssignor.java}: \method{Map<String,List<TopicPartition>> assign(List<TopicPartition> partitions, List<String> members)}. Implement only \method{assign}.}
\subsubsection*{Inputs, outputs, and state}
\begin{itemize}
\item \pathcode{partitions} is the list of \pathcode{TopicPartition} values to assign; \pathcode{members} is the list of member identifiers. Neither list nor any element may be null. Keep only one copy of duplicate partitions or members.
\item Sort partitions by topic then partition number, and members by Java \pathcode{String} natural order; do not preserve input order or sort numeric-looking strings by numeric value.
\item The returned map contains every distinct member, including members with empty lists; both the map and its lists are unmodifiable. The assignor does not change inputs and has no state between calls.
\item Null inputs/elements, or a member identifier outside 1--255 UTF-8 bytes, report \method{INVALID_REQUEST}. With no members, return an empty map; with members but no partitions, retain every member with an empty list.
\end{itemize}
\lessoncontract{When there is at least one distinct member, every distinct partition appears exactly once and is owned by a member; every distinct member appears in the result. The maximum and minimum partition counts among members differ by no more than one. With no members, return the empty map as stated above; no partition has an owner.}
\expected{Start with the explicit partition list \pathcode{[orders-0, orders-1, orders-2, orders-3, orders-4]} and members in input order \pathcode{[b,a,a]}. After deduplication and sorting, members are [a,b]; a gets 0, 2, 4 and b gets 1, 3. Shuffling input does not change the result. With five partitions and members sorted by Java string order as \pathcode{[a,b,c,d,e,f,g]}, the first five get one each and f/g remain with empty lists.\\
\method{Step17Test.sortsAndDeduplicatesPartitionsAndMembersBeforeRoundRobinAssignment} checks deduplication and deterministic order; \method{assignsFivePartitionsToTheFirstFiveOfSevenSortedMembersAndKeepsEmptyMembers} checks empty assignments; \method{sortsCrossTopicPartitionsAndMemberIdsLexicallyAfterDeduplication} checks cross-topic and lexical order, where \pathcode{"10"} precedes \pathcode{"2"}.}
\lessonhints{Break “one owner” into two independently checkable properties: complete coverage and no duplicates.}{Sort members and partitions first, then assign by index modulo member count; the input List order is not the course rule.}{Work out five partitions with two members and five partitions with seven members, then add duplicate inputs and verify no partition appears twice.}
\lessonquestions{Why is stable sorting observable contract rather than output polish?}{When there are more members than partitions, why return the members with empty lists?}
\paragraph{Self-check explanation.}Stable sorting gives the same owner for the same sets on every call, which makes coordination and verification reproducible; an empty list says a member remains in the group but currently owns no partition.
\pitfalls{Returning only members that own partitions loses explicit empty assignments; depending on incidental map iteration breaks reproducibility. Real Kafka has several assignors and cooperative rebalancing; this course implements only a simple single-topic Round Robin and does not claim it is Kafka's default algorithm.}
\lessoncommand{17}

\lessonsection{18}{Membership and Rebalancing}
\goals{Manage group members, generation, and the full assignment, rebalancing correctly when membership changes.}{Understand what join, leave, generation, and committed offset each represent.}
\background{A member is a consumer instance identified by a member id; in this course a group subscribes to one topic. The coordinator stores the member set and its shared generation, and uses the Step 17 assignor to calculate the complete assignment. A new member joining or an existing member leaving changes the topology, increments the generation, and immediately recomputes partitions for every member. Repeating join for the same member and topic refreshes only its heartbeat; it does not change topology. A generation is neither a timestamp nor consumer progress.}
\subsubsection*{Read first}
Review \stepref{10}{topic metadata and the partition catalog}, then \stepref{17}{deterministic assignment}. You need to know which \pathcode{TopicPartition} values belong to a topic and distinguish the coordinator's in-memory membership from the durable offsets in Step 15.
\providededit{\method{PartitionCatalog}, group lock, member-state container, \method{TimeSource}, and the Step 17 assignor are provided; message request value types are already encoded.}{\pathcode{exercises/src/main/java/io/simplekafka/group/GroupCoordinator.java}: \method{GroupAssignment join(String group, String member, String topic)}, \method{GroupAssignment leave(String group, String member)}, and \method{GroupAssignment assignment(String group, String member)}; extend \method{handle(Messages.Request)} in \pathcode{exercises/src/main/java/io/simplekafka/broker/BrokerHandler.java} with \pathcode{JOIN_GROUP/LEAVE_GROUP/GROUP_ASSIGNMENT}.}
\subsubsection*{Inputs, outputs, and state}
\begin{itemize}
\item group/member identifiers must be nonempty and at most 255 UTF-8 bytes; topic must be a valid course topic. A first join creates the group only after verifying that the topic exists, then advances its initial generation from 0 to 1.
\item A new member joining or a known member successfully leaving increments generation exactly once and returns a \method{GroupAssignment} for all current members. Repeating join for the same member and topic updates only lastHeartbeat and returns the existing assignment; it must not duplicate the member or increment generation.
\item A topic conflict in an existing group takes priority and reports \method{INVALID_REQUEST}; an unknown topic reports \method{UNKNOWN_TOPIC_OR_PARTITION} and does not create a group or consume a generation. Invalid identifiers report \method{INVALID_REQUEST}. leave or assignment for an unknown group/member reports \method{UNKNOWN_MEMBER}; operations after coordinator close throw \pathcode{IllegalStateException}.
\item When the group becomes empty, retain its generation history in the coordinator and retain committed offsets; offset storage is separate from member lifetime. Membership itself is process-local and does not become a durable group protocol just because offsets persist.
\end{itemize}
\lessoncontract{Each real topology change increments one generation and recalculates assignments for every remaining member; a repeated join does not change the generation. Assignment queries succeed only for a currently known member and return the same full-group mapping.}
\expected{Initially \pathcode{workers} does not exist, and \pathcode{orders} has partitions 0, 1, and 2. a joins and receives generation 1 with all three partitions; a joins the same topic again and generation remains 1. b joins, generation 2, with a assigned 0 and 2 and b assigned 1. c joins at generation 3 and each member receives one partition. c leaves at generation 4; the new assignment is a: 0,2 and b: 1. If the group later becomes empty, its generation remains and its existing group offsets can still be fetched.\\
\method{Step18Test.membershipChangesRebalanceEverySurvivorAndRetainEmptyGroupGenerationOverTcp} uses real TCP to check join, repeated join, full-group rebalance, topic/unknown-member errors, empty-group offset retention, and a later rejoin.}
\lessonhints{First decide whether a request actually changes the member set; a repeated join, a rejected join, and a new member join are different state transitions.}{After each topology change, check that every member sees the same generation and full assignment, not only that the joiner was updated.}{Starting from no members, list the generation and every owner after a joins, a repeats, b joins, c joins, and c leaves.}
\lessonquestions{Why must a receive an assignment for a new generation when b joins?}{Why should group offsets not be deleted when the last member leaves?}
\paragraph{Self-check explanation.}A partition may have moved from a to b, so a's old assignment no longer agrees with server ownership; consumer progress belongs to a group/topic-partition, not to an online member, so an empty group can retain it.
\pitfalls{Advancing a separate version per member makes tokens disagree; an empty group does not mean its consumer progress is invalid. Real Kafka uses multi-phase join, assignment synchronization, cooperative revocation, and other protocol steps; the course rebalances immediately and does not simulate the full classic barrier or KIP-848.}
\lessoncommand{18}

\lessonsection{19}{Heartbeats and Expiration}
\goals{Detect idle members with a controllable clock and process all same-group expirations as one topology change.}{Understand lastHeartbeat, session timeout, and the inclusive equality boundary.}
\background{A member process may crash without leaving, so the coordinator cannot immediately remove it just because it is silent. A successful heartbeat records the current millisecond value from the injected clock and indicates that the member is active. A member expires once elapsed time reaches the session timeout. A deterministic clock makes boundary tests independent of sleep; the provided periodic worker is the runtime trigger, while a test's direct expiry hook is a separate mechanism.}
\subsubsection*{Read first}
You need \stepref{18}{members, generation, and assignment}. Be able to explain that members in a group share a generation and that a removed member's old token is no longer a valid heartbeat.
\providededit{The \method{TimeSource}, group lock, positive session timeout, and periodic scheduling lifecycle are provided.}{\pathcode{exercises/src/main/java/io/simplekafka/group/GroupCoordinator.java}: \method{void heartbeat(GroupToken token)}, \method{Set<String> expire()}; extend \method{handle(Messages.Request)} in \pathcode{exercises/src/main/java/io/simplekafka/broker/BrokerHandler.java} with \pathcode{HEARTBEAT}.}
\subsubsection*{Inputs, outputs, and state}
\begin{itemize}
\item The coordinator's sessionTimeoutMillis must be positive at construction or reports \method{INVALID_REQUEST}. A null heartbeat token reports \method{INVALID_REQUEST}; check member existence first, reporting \method{UNKNOWN_MEMBER} when absent; then check generation, reporting \method{ILLEGAL_GENERATION} for a stale or future value. Failures do not refresh time; only a valid token updates that member's lastHeartbeat to the current clock. Operations after coordinator close throw \pathcode{IllegalStateException}.
\item \method{expire()} takes no argument and uses the injected current clock. It returns the sorted set of removed-member identifiers formatted as \pathcode{group/member}. Expiration uses \pathcode{now-lastHeartbeat >= sessionTimeoutMillis}; a member below the threshold remains, and a future heartbeat timestamp does not expire early.
\item One expire call can remove several expired members from one group, but increments that group's generation once and recalculates assignment once; non-expired members remain. An empty result means no member was removed; repeating expiry makes no further change. Committed offsets are not deleted.
\item The broker's periodic worker triggers checks automatically. For deterministic tests, move the \method{TimeSource} and call the provided \pathcode{expireGroups()} hook or \method{expire()}, rather than guessing timing with real sleep.
\end{itemize}
\lessoncontract{The expiry boundary is \pathcode{>=}; an invalid/stale heartbeat does not renew membership. Results sort by group and then member. Each group changed by one expiry pass increments generation once, and an empty group retains its generation.}
\expected{Set clock=1000 and timeout=100; a and b join \pathcode{workers} at 1000, making generation=2, so a's generation-1 heartbeat is already stale. At 1099, a's old token and future generation 99 report \method{ILLEGAL_GENERATION}, and an unknown member reports \method{UNKNOWN_MEMBER}; none refreshes time. b's valid generation-2 heartbeat updates its time to 1099. expire at 1099 returns empty; at 1100 it removes only \pathcode{workers/a}, exactly 100 ms after a's last heartbeat. Generation becomes 3 and b receives all three partitions.\\
\method{Step19Test.rejectsStaleAndFutureHeartbeatsAndExpiresOnlyTheUnrefreshedMemberAtTheBoundary} checks 1099/1100 and failed heartbeats; \method{expirationOrdersAllIdentifiersAndAdvancesEachChangedGroupOnce} checks sorting and one increment per group; \method{idempotentJoinRefreshesHeartbeatButRejectedJoinDoesNot} checks renewal; \method{replicatedClusterExpiresIdleMembersWithoutIncomingRequests} and \method{expireGroupsProvidesDeterministicClusterExpiryAndChecksLifecycle} check the periodic worker and deterministic hook.}
\lessonhints{Use the current clock, check elapsed time, and remove members within one consistent state observation; a failed heartbeat must not write lastHeartbeat.}{The return value should identify the exact group/member, not only the number of expirations.}{With timeout=100 and lastHeartbeat=1000 compare 1099 and 1100; refresh B at 1099 and confirm that this pass removes only A.}
\lessonquestions{Why must the equality boundary be explicit in the timeout contract?}{If A and C in one group expire in the same call, why not increment generation twice?}
\paragraph{Self-check explanation.}Without an equality rule, ownership at the exact deadline would be ambiguous; one expiry pass observes one topology change for that group, so removing its expired members together triggers one full rebalance.
\pitfalls{Using system time and sleep for an exact threshold is scheduler-sensitive; use an injected clock for boundaries and test the worker separately as integration behavior. Updating lastHeartbeat before validating a token lets an expired member renew illegally. Real Kafka heartbeat behavior also depends on configuration, network, and coordination protocol.}
\lessoncommand{19}

\lessonsection{20}{Consumer Group Fencing}
\goals{Combine assignments with fetch and offset commit, using generation tokens to reject stale operations from zombie members.}{Understand the token triple, revoked partitions, assignment refresh, and the boundary around external side effects.}
\background{An old consumer thread may still run after a rebalance. A GroupToken=(group,member,generation) is the group-membership version carried by a request; before allowing tokenized FETCH or COMMIT, the broker must confirm that the member exists, the generation is current, and the member owns the target partition. Fencing can reject stale requests, but it cannot recall records already delivered to the old client or undo business side effects already performed. Manual mode has no token and continues to use the earlier ordinary consumer path.}
\subsubsection*{Read first}
You need \stepref{14}{fetch positions}, \stepref{15}{durable offset keys}, \stepref{17}{assignment algorithm}, \stepref{18}{generation and assignment}, and \stepref{19}{heartbeat and expiry}. This step combines them; do not treat a member generation as a consumer offset.
\providededit{Group request/response types, \method{RpcClient}, the \method{SimpleConsumer} pull/offset path, and member-lifecycle skeleton are provided.}{\pathcode{exercises/src/main/java/io/simplekafka/group/GroupCoordinator.java}: \method{void validate(GroupToken token, TopicPartition tp)}; \pathcode{exercises/src/main/java/io/simplekafka/client/GroupConsumer.java}: \method{GroupConsumer(RpcClient client, String group, String member)} (takes ownership and closes the client), \method{GroupConsumer(MetadataRouter router, String group, String member)} (borrows the router and does not close it); both constructors require group/member identifiers of 1--255 UTF-8 bytes. Implement only \method{subscribe(String topic)}, \method{Map<TopicPartition,List<LogRecord>> poll(int maxRecords, int maxBytes)}, and \method{commitSync()}; extend \pathcode{exercises/src/main/java/io/simplekafka/broker/BrokerHandler.java} with token fencing for \pathcode{FETCH/COMMIT}.}
\subsubsection*{Inputs, outputs, and state}
\begin{itemize}
\item \method{GroupCoordinator.validate(token,tp)} requires a current member, current generation, and ownership of the target partition. Unknown member reports \method{UNKNOWN_MEMBER}; a mismatched generation, including a future token, reports \method{ILLEGAL_GENERATION}; a partition not assigned to the member reports \method{NOT_ASSIGNED}. A null tp reports \method{INVALID_REQUEST}; a rejected validation does not write OffsetStore.
\item \method{subscribe(topic)} explicitly joins and obtains an assignment/token. A newly assigned partition starts at its committed offset or 0; a retained partition keeps its position; a revoked partition is dropped. Poll/commit before subscribing reports \method{INVALID_REQUEST}; unknown-topic and other errors propagate. Operations after consumer close throw \pathcode{IllegalStateException}.
\item \method{poll(maxRecords,maxBytes)} requires positive total budgets or reports \method{INVALID_REQUEST}, and heartbeats before fetch requests; its returned map is unmodifiable. If the heartbeat reports \method{ILLEGAL_GENERATION}, it only queries the current assignment, updates the local token/positions, then fetches with the new token: revoked positions are removed and new partitions start at committed/0. On \method{UNKNOWN_MEMBER}, it does not rejoin automatically; the caller must explicitly subscribe. Other errors propagate. A successful fetch advances the relevant in-memory position but does not itself commit.
\item \method{commitSync()} sends the current token with each assigned partition's position. It neither heartbeats nor refreshes the assignment. A stale token/non-owner rejection leaves the store unchanged; the per-partition requests are not atomic, so an earlier commit may succeed before a later request fails. The broker COMMIT branch requires a nonnegative nextOffset, and the OffsetKey's group and partition must match the token and request; otherwise it reports \method{INVALID_REQUEST} without append. The server validates and appends under the coordinator lock, with lock order coordinator then store; this is not a partition-log transaction. A null token remains valid on the ordinary manual requests.
\item Successful token FETCH returns \pathcode{FetchBody}, token COMMIT returns \pathcode{EmptyBody}, and assignment lookup returns \pathcode{GroupBody}. Previously sent records and client side effects are not rolled back by a later rejection. Operations after coordinator close throw \pathcode{IllegalStateException}.
\end{itemize}
\lessoncontract{Both FETCH and COMMIT fence member, generation, and partition owner. The COMMIT OffsetKey must match the token's group and the requested partition; any rejection must not append that offset. After a generation changes, an old member cannot commit a revoked partition's cached position.}
\expected{Initially \pathcode{orders-0} contains only offset 0; \pathcode{orders-1} contains offsets 0--8, and \pathcode{workers/orders-1} is committed at 7. a subscribes at generation 1 and owns both partitions; poll(1,4096) reads p0/offset0, leaving positions p0=1 and p1=7. b subscribes, causing generation 2; a now owns only p0, and b owns p1. a's old-token commitSync reports \method{ILLEGAL_GENERATION}; p0 is still uncommitted and p1 is still 7. Append p0/offset1. a polls, detects the old generation, obtains the new assignment, retains p0=1, drops revoked p1, and reads p0/offset1. Its commit stores p0=2 while p1 remains 7. b's first poll starts at committed p1=7. If the broker fixture restarts, in-memory membership disappears but the log and offsets remain (p0=2, p1=9). Create a new \method{GroupConsumer} a and explicitly subscribe; the new generation starts at 1, and its positions resume from committed offsets 2 and 9.\\
\method{Step20Test.coordinatorFencesCommitsByGenerationOwnershipAndMatchingOffsetKey} checks the coordinator and mismatched keys; \method{rawTcpFetchAndCommitTokensFenceInvalidOwnersWithoutChangingOffsets} checks stale/future/unknown/non-owner tokens over real TCP; \method{groupConsumerPreservesRetainedPositionDropsRevokedPartitionsAndRecoversAfterExpiryAndRestart} checks refresh, takeover, and restart recovery; \method{groupConsumerClosesItsOwnedRpcClientAfterPollingARecord} checks resource ownership.}
\lessonhints{Request validation must check member, generation, and partition owner together; comparing only the group name is not fencing.}{After rebalancing, classify each partition as retained with a local position, newly acquired and restored from committed/0, or revoked and discarded.}{Use the example to have a submit a larger offset with its old generation; check the error code and verify every related OffsetKey remains unchanged.}
\lessonquestions{What can generation fencing reject, and why can it not guarantee exactly-once external side effects?}{If validation unlocks before commit, what might happen between the two operations?}
\paragraph{Self-check explanation.}The server can reject a stale token's fetch/commit, but cannot recall a delivered record or undo an action in an external system. If a rebalance occurs between validation and append, a zombie could write progress after losing ownership.
\pitfalls{Validating, releasing the lock, and committing later leaves a rebalance race. Turning both \method{ILLEGAL_GENERATION} and \method{UNKNOWN_MEMBER} into automatic rejoin hides different states. Real Kafka group protocols and transactional commits are more complex; this course's in-process coordinator and local offset store are not a highly available \pathcode{__consumer_offsets}.}
\lessoncommand{20}
